\chapter{Lasers \& Optical Resonance}

\section{Introduction and Historical Evolution}
The concept of stimulated emission was first proposed theoretically by Albert Einstein in 1917. However, it took nearly four decades before Charles Townes built the microwave MASER (1953) and Theodore Maiman constructed the first operational optical LASER in 1960 using a synthetic ruby crystal. 

Today, lasers span an enormous range of engineering applications, from microscopic semiconductor diodes powering global fiber communications and barcode scanners to multi-kilowatt industrial cutting lasers, medical surgical scalpels, and multi-terawatt inertial confinement fusion systems.

\section{Fundamental Laser Characteristics}
Laser light differs fundamentally from conventional thermal radiation across four core dimensions:
\begin{enumerate}[leftmargin=*,itemsep=3pt]
    \item \textbf{Monochromaticity:} Thermal light sources emit broad continuous spectrums or broad spontaneous lines ($\Delta\lambda \sim 10\text{--}100\text{ nm}$). In contrast, laser emission is confined to an extremely narrow spectral linewidth ($\Delta\nu \sim 1\text{ kHz}$ to $1\text{ MHz}$, corresponding to $\Delta\lambda \ll 10^{-4}\text{ nm}$).
    \item \textbf{Coherence:}
    \begin{itemize}
        \item \textbf{Temporal Coherence:} Quantifies the interval over which the light wave maintains a predictable, uninterrupted phase. The coherence length is $L_c = c \tau_c \approx \frac{c}{\Delta\nu} = \frac{\lambda^2}{\Delta\lambda}$. For He-Ne lasers, $L_c$ can exceed several kilometers.
        \item \textbf{Spatial Coherence:} Quantifies the phase correlation across different points on the wavefront perpendicular to the direction of propagation. This enables lasers to be focused onto diffraction-limited spots of size $d \approx \lambda$.
    \end{itemize}
    \item \textbf{Directionality:} Lasers emit nearly perfect parallel planar wavefronts. The angular divergence $\theta$ is governed strictly by wave diffraction: $\theta \approx 1.22 \frac{\lambda}{D}$, where $D$ is the cavity aperture diameter. Typical beam divergences are under $1\text{ milliradian}$.
    \item \textbf{High Brightness / Intensity:} Because enormous optical energy is concentrated into a tight spatial cross-section and narrow bandwidth, laser surface brightness exceeds the surface brightness of the Sun by orders of magnitude.
\end{enumerate}

\section{Quantum Interaction of Radiation with Matter}
Consider an idealized atomic system with discrete ground state $E_1$ and excited state $E_2$ separated by energy transition gap $\Delta E = E_2 - E_1 = h\nu$.

\subsection{1. Induced (Stimulated) Absorption}
When an external photon of energy $h\nu$ collides with an atom in the lower energy level $E_1$, the photon is annihilated, and the atom transitions to the upper level $E_2$. The number of absorption transitions per unit volume per unit time is:
\begin{equation}
R_{12} = B_{12} N_1 u(\nu)
\end{equation}
where $N_1$ is the ground-state population density ($\text{atoms/m}^3$), $u(\nu)$ is the spectral radiation energy density ($\text{J}\cdot\text{s/m}^3$), and $B_{12}$ is Einstein's coefficient of stimulated absorption.

\subsection{2. Spontaneous Emission}
An atom residing in excited state $E_2$ is inherently unstable. After a characteristic spontaneous lifetime $\tau_{\text{sp}} \sim 10^{-8}\text{ s}$, it transitions down to $E_1$, emitting a photon of energy $h\nu$. This process occurs independently without external electromagnetic trigger. The emitted photon has a completely random propagation direction, random polarization, and random phase. The rate is:
\begin{equation}
R_{21}^{\text{spont}} = A_{21} N_2
\end{equation}
where $A_{21}$ is Einstein's coefficient of spontaneous emission.

\subsection{3. Stimulated Emission}
If an incoming photon of exact resonant energy $h\nu = E_2 - E_1$ interacts with an atom already residing in excited state $E_2$ before it decays spontaneously, the external electric field forces the electron to drop to $E_1$. The atom emits a second photon that is an \textbf{exact coherent clone} of the incident trigger photon, possessing:
\begin{itemize}[itemsep=1pt]
    \item Exactly identical frequency $\nu$.
    \item Identical phase ($\Delta\phi = 0$).
    \item Identical propagation direction vector $\vec{k}$.
    \item Identical polarization orientation $\hat{e}$.
\end{itemize}
The stimulated transition rate is:
\begin{equation}
R_{21}^{\text{stim}} = B_{21} N_2 u(\nu)
\end{equation}
where $B_{21}$ is Einstein's coefficient of stimulated emission.

\section{Rigorous Derivation of Einstein's A and B Relations}
Consider an atomic assembly enclosed inside an isothermal cavity at absolute temperature $T$, in thermodynamic equilibrium with blackbody radiation. Under steady-state equilibrium:
\begin{equation}
\text{Total Upward Transitions} = \text{Total Downward Transitions}
\end{equation}
\begin{equation}
R_{12} = R_{21}^{\text{spont}} + R_{21}^{\text{stim}} \implies B_{12} N_1 u(\nu) = A_{21} N_2 + B_{21} N_2 u(\nu)
\end{equation}
Rearranging terms to isolate the radiation energy density $u(\nu)$:
\begin{equation}
u(\nu) [B_{12} N_1 - B_{21} N_2] = A_{21} N_2 \implies u(\nu) = \frac{A_{21} N_2}{B_{12} N_1 - B_{21} N_2} = \frac{A_{21}/B_{21}}{\left(\frac{B_{12}}{B_{21}}\frac{N_1}{N_2}\right) - 1}
\end{equation}
According to Maxwell-Boltzmann statistical mechanics, the population ratio of states in thermodynamic equilibrium is:
\begin{equation}
\frac{N_1}{N_2} = \frac{g_1}{g_2} \exp\left(\frac{E_2 - E_1}{k_B T}\right) = \frac{g_1}{g_2} e^{h\nu / k_B T}
\end{equation}
where $g_1$ and $g_2$ are the statistical degeneracies of levels 1 and 2. Substituting into the energy density expression:
\begin{equation}
u(\nu) = \frac{A_{21}/B_{21}}{\left(\frac{B_{12}}{B_{21}}\frac{g_1}{g_2}\right) e^{h\nu/k_B T} - 1}
\end{equation}
According to Planck's Law of Blackbody Radiation:
\begin{equation}
u(\nu) = \frac{8\pi h \nu^3}{c^3}\left( \frac{1}{e^{h\nu/k_B T} - 1} \right)
\end{equation}
For these two fundamental equations to be identical across all temperatures $T$ and frequencies $\nu$, the corresponding coefficients must match:
\begin{align}
\frac{B_{12}}{B_{21}}\frac{g_1}{g_2} &= 1 \implies B_{12} = \frac{g_2}{g_1} B_{21} \quad (\text{For non-degenerate levels, } B_{12} = B_{21}) \\
\frac{A_{21}}{B_{21}} &= \frac{8\pi h \nu^3}{c^3} = \frac{8\pi h}{\lambda^3}
\end{align}

\begin{conceptbox}{Physical Consequence of the $\nu^3$ Scaling Law}
The ratio of spontaneous emission to stimulated emission rate is:
\begin{equation}
\frac{R_{\text{spont}}}{R_{\text{stim}}} = \frac{A_{21}}{B_{21} u(\nu)} = e^{h\nu/k_B T} - 1
\end{equation}
\begin{itemize}[itemsep=2pt]
    \item In the \textbf{Microwave Regime} ($\nu \approx 10^{10}\text{ Hz}$): $h\nu \ll k_B T \implies e^{h\nu/k_B T} - 1 \approx \frac{h\nu}{k_B T} \ll 1$. Stimulated emission easily dominates, allowing masers to operate without intense optical cavities.
    \item In the \textbf{Visible/Optical Regime} ($\nu \approx 5 \times 10^{14}\text{ Hz}, \lambda \approx 600\text{ nm}$): At room temperature ($T = 300\text{ K}$), $h\nu \approx 2.07\text{ eV} \gg k_B T \approx 0.026\text{ eV}$.
    \[
    e^{h\nu/k_B T} \approx e^{80} \approx 5.5 \times 10^{34} \gg 1
    \]
    Spontaneous emission dominates completely by 34 orders of magnitude under thermal equilibrium! Thus, coherent laser action is physically impossible without:
    \begin{enumerate}
        \item Strong non-equilibrium external pumping to create \textbf{population inversion}.
        \item An optical resonator cavity to establish massive photon field densities $u(\nu)$.
    \end{enumerate}
\end{itemize}
\end{conceptbox}

\section{Essential Architectural Elements of a Laser}
Every operational laser consists of three fundamental subsystems:
\begin{enumerate}[leftmargin=*,itemsep=3pt]
    \item \textbf{Active Lasing Medium:} The physical collection of atoms, gas molecules, semiconductor crystals, or doped dielectric glass containing suitable quantum transitions with a long-lived \textbf{metastable state} ($\tau_{\text{meta}} \sim 10^{-3}\text{ s}$ compared to normal $\tau \sim 10^{-8}\text{ s}$).
    \item \textbf{Pumping System:} Provides external excitation energy to populate the upper state. Primary methods:
    \begin{itemize}
        \item \textbf{Optical Pumping:} Intense xenon flashlamps or diode laser arrays (used in solid-state lasers like Ruby and Nd:YAG).
        \item \textbf{Electrical Gas Discharge:} High-voltage electron impact excitation (used in He-Ne, $\text{CO}_2$, and Argon-ion lasers).
        \item \textbf{Direct Injection Current:} Forward biasing a semiconductor p-n junction (diode lasers).
        \item \textbf{Chemical Reaction:} Exothermic gas reactions creating inverted states (HF/DF chemical lasers).
    \end{itemize}
    \item \textbf{Optical Resonant Cavity (Fabry-Pérot Resonator):} A pair of highly polished mirrors facing each other along the active medium axis:
    \begin{itemize}
        \item Rear Mirror: 100\% high-reflectivity mirror.
        \item Output Coupler: Partially transmissive mirror ($R \approx 95\text{--}99\%$) through which the useful output beam emerges.
    \end{itemize}
\end{enumerate}

\section{Population Inversion, Gain Coefficient, and Threshold Condition}
When a collimated monochromatic beam passes through a medium of length $dz$, the change in intensity is:
\begin{equation}
dI = I(z) \sigma_{21} \left( N_2 - \frac{g_2}{g_1} N_1 \right) dz \equiv \gamma(\nu) I(z) dz
\end{equation}
where $\sigma_{21}$ is the stimulated emission cross-section and $\gamma(\nu) = \sigma_{21}(N_2 - N_1)$ is the \textbf{small-signal gain coefficient}.
Integrating over path length $z$:
\begin{equation}
I(z) = I_0 e^{\gamma(\nu) z}
\end{equation}
If $N_2 < N_1$, $\gamma(\nu) < 0$ and the beam suffers Beer-Lambert exponential absorption. If $N_2 > N_1$ (\textbf{Population Inversion}), $\gamma(\nu) > 0$ and the beam undergoes exponential optical amplification.

\subsection{Cavity Threshold Equation}
During one complete round trip of length $2L$ in a cavity with mirror reflectivities $R_1, R_2$ and distributed internal scattering loss $\alpha_s$, the round-trip gain condition for sustained oscillation is:
\begin{equation}
R_1 R_2 e^{2(\gamma_{\text{th}} - \alpha_s)L} = 1
\end{equation}
Taking the natural logarithm yields the \textbf{Threshold Gain}:
\begin{equation}
\gamma_{\text{th}} = \alpha_s + \frac{1}{2L}\ln\left(\frac{1}{R_1 R_2}\right)
\end{equation}

\section{Three-Level vs Four-Level Laser Schemes}
\begin{itemize}[leftmargin=*,itemsep=3pt]
    \item \textbf{Three-Level System (e.g., Ruby Laser):} Ground state $E_1$ acts as both the initial state and the lower lasing terminal level. To achieve population inversion ($N_2 > N_1$), more than 50\% of all atoms in the active crystal must be pumped into the excited state simultaneously. This requires enormous pump power, making continuous-wave (CW) operation extremely difficult (operates in giant pulses).
    \item \textbf{Four-Level System (e.g., He-Ne, Nd:YAG):} Lasing transitions terminate on an intermediate level $E_2$ located well above the ground state $E_1$ ($E_2 - E_1 \gg k_B T$). Because $E_2$ is thermally depopulated ($N_2 \approx 0$) and decays rapidly to $E_1$ ($\tau_{21} \sim 10^{-11}\text{ s}$), even a modest population in metastable level $E_3$ easily achieves $N_3 > N_2$. Four-level lasers have dramatically lower threshold pumping powers and operate continuously in CW mode with high stability.
\end{itemize}

\section{Major Practical Engineering Lasers}

\subsection{1. Ruby Laser (3-Level Solid State)}
\begin{itemize}[itemsep=2pt]
    \item \textbf{Active Medium:} Single-crystal corundum ($\text{Al}_2\text{O}_3$) doped with $0.05\%$ by weight of Chromium ions ($\text{Cr}^{3+}$).
    \item \textbf{Pumping:} Helical Xenon flashlamp pulsing at milliseconds duration.
    \item \textbf{Lasing Transition:} Non-radiative transition from $^4F_2, ^4F_1$ pump bands populates the $^2E$ metastable doublet state ($\tau \sim 3\text{ ms}$). Radiative laser emission occurs to ground state $^4A_2$ at $\lambda = 694.3\text{ nm}$ (deep red).
\end{itemize}

\subsection{2. Helium-Neon (He-Ne) Gas Laser (4-Level Gas)}
\begin{itemize}[itemsep=2pt]
    \item \textbf{Active Medium:} Gas mixture of Helium and Neon in a 10:1 ratio at total pressure $\approx 1\text{ torr}$.
    \item \textbf{Pumping:} DC electrical discharge excites He atoms via electron impact to metastable states $2^3S_1$ ($19.82\text{ eV}$) and $2^1S_0$ ($20.61\text{ eV}$).
    \item \textbf{Resonant Energy Transfer:} Excited He atoms collide inelastically with ground-state Ne atoms, transferring energy with high efficiency to closely matched Ne $2s$ and $3s$ levels ($20.66\text{ eV}$).
    \item \textbf{Laser Emission:} The primary transition $3s_2 \to 2p_4$ yields the standard red line at $\lambda = 632.8\text{ nm}$. Subsequent fast spontaneous decay from $2p_4 \to 1s$ preserves continuous population inversion.
\end{itemize}

\subsection{3. Nd:YAG Laser (4-Level Solid State)}
\begin{itemize}[itemsep=2pt]
    \item \textbf{Active Medium:} Yttrium Aluminum Garnet ($\text{Y}_3\text{Al}_5\text{O}_{12}$) doped with $1\%\text{ Nd}^{3+}$.
    \item \textbf{Emission:} Pumping at $808\text{ nm}$ creates inversion in the $^4F_{3/2}$ metastable level. Transition to $^4I_{11/2}$ produces high-power near-infrared emission at $\lambda = 1064\text{ nm}$. Frequency-doubling with KDP/BBO non-linear crystals yields green light at $532\text{ nm}$.
\end{itemize}

\subsection{4. Semiconductor Injection Laser (GaAs)}
\begin{itemize}[itemsep=2pt]
    \item \textbf{Active Medium:} Heavily doped p-n junction of Gallium Arsenide (GaAs).
    \item \textbf{Pumping:} Forward bias injects high-density electrons and holes into the narrow active depletion region.
    \item \textbf{Direct Bandgap Recombination:} Radiative band-to-band carrier recombination generates stimulated photons: $\lambda = \frac{hc}{E_g} \approx \frac{1240\text{ eV}\cdot\text{nm}}{1.424\text{ eV}} \approx 870\text{ nm}$.
\end{itemize}


\section{Optical Resonator Modes and Beam Properties}
An optical resonator consists of two mirrors with radii of curvature $R_1, R_2$ separated by cavity length $L$.
\subsection{Cavity Stability Criterion}
Defining cavity $g$-parameters: $g_1 = 1 - \frac{L}{R_1}$ and $g_2 = 1 - \frac{L}{R_2}$. The condition for stable ray confinement without lateral leakage is:
\begin{equation}
0 \le g_1 g_2 \le 1
\end{equation}

\subsection{Transverse Electromagnetic Modes (TEM)}
In addition to longitudinal modes $\nu_m = m \frac{c}{2L}$, the field distribution across the mirror cross-section forms discrete transverse patterns $\text{TEM}_{mn}$:
\begin{itemize}[itemsep=2pt]
    \item \textbf{Fundamental Mode ($\text{TEM}_{00}$):} Gaussian intensity profile $I(r) = I_0 \exp\left(-\frac{2r^2}{w(z)^2}\right)$ with zero nulls. It possesses minimum possible beam divergence and maximum focusability.
    \item \textbf{Higher-Order Modes ($\text{TEM}_{01}, \text{TEM}_{11}$):} Exhibit nodal lines with Hermite-Gaussian field distributions.
\end{itemize}
The beam waist radius expands with propagation distance $z$ according to:
\begin{equation}
w(z) = w_0 \sqrt{1 + \left(\frac{z}{z_R}\right)^2}, \quad \text{where } z_R = \frac{\pi w_0^2}{\lambda} \text{ is the Rayleigh Range}
\end{equation}

\section{Giant Pulse Techniques: Q-Switching and Mode-Locking}
\begin{enumerate}[leftmargin=*,itemsep=3pt]
    \item \textbf{Q-Switching (Quality Factor Modulation):}
    \begin{itemize}
        \item While pumping, cavity losses are kept artificially high (low $Q$-factor), suppressing laser action and allowing massive energy storage in the metastable level ($N_2 \gg N_{\text{th}}$).
        \item The cavity loss is abruptly switched to zero (high $Q$-factor), releasing all stored population inversion in a single giant nanosecond pulse ($P_{\text{peak}} \sim 10^6\text{--}10^9\text{ W}$, pulse duration $\sim 10\text{--}50\text{ ns}$).
        \item Methods: Electro-optic Pockels cells, acousto-optic modulators, saturable dye absorbers.
    \end{itemize}
    \item \textbf{Mode-Locking:}
    \begin{itemize}
        \item Longitudinal cavity modes are locked in a fixed phase relationship ($\Delta\phi = 0$).
        \item Constructive interference produces periodic ultra-short femtosecond pulses ($\tau_p \sim 10\text{--}100\text{ fs}$) at repetition rate $f_{\text{rep}} = \frac{c}{2L}$.
    \end{itemize}
\end{enumerate}

\section{Principles of Holography}
Conventional 2D photography records only the square of the electric field amplitude ($I \propto |E|^2$), completely discarding wavefront phase information $\phi(x,y)$. Dennis Gabor (Nobel Prize 1971) invented \textbf{Holography} to record both amplitude and phase:
\begin{enumerate}[leftmargin=*,itemsep=2pt]
    \item \textbf{Recording Stage:} A coherent laser beam is split into a reference wave $R(x,y) = A_r e^{i\phi_r}$ and an illumination wave that scatters off the 3D object to form object wave $O(x,y) = A_o e^{i\phi_o}$. The interference pattern recorded on high-resolution emulsion is:
    \begin{equation}
    I(x,y) = |R + O|^2 = |R|^2 + |O|^2 + R^* O + R O^*
    \end{equation}
    \item \textbf{Reconstruction Stage:} When illuminated by the original reference wave $R$:
    \begin{equation}
    \Psi_{\text{out}} = R \cdot I(x,y) = R(|R|^2 + |O|^2) + |R|^2 O + R^2 O^*
    \end{equation}
    The term $|R|^2 O = A_r^2 A_o e^{i\phi_o}$ is an exact spatial duplicate of the original object wavefront, creating a true 3D virtual image with full parallax depth perception.
\end{enumerate}

\section{Worked Numerical Examples}

\begin{examplebox}{Worked Example 2.1: Cavity Mode Spacing and Population Ratio}
\textbf{Problem:} A He-Ne laser has a cavity length $L = 30\text{ cm}$ operating at $\lambda = 632.8\text{ nm}$. (a) Calculate the frequency spacing $\Delta\nu$ between adjacent longitudinal cavity modes. (b) Find the ratio of spontaneous to stimulated emission rates at room temperature ($T = 300\text{ K}$).\\
\textbf{Solution:}\\
(a) The longitudinal mode condition is $L = m \frac{\lambda}{2} \implies \nu_m = m \frac{c}{2L}$. The inter-mode frequency separation is:
\[
\Delta\nu = \nu_{m+1} - \nu_m = \frac{c}{2L} = \frac{3 \times 10^8\text{ m/s}}{2 \times 0.30\text{ m}} = 5 \times 10^8\text{ Hz} = 500\text{ MHz}
\]
(b) The photon energy is:
\[
h\nu = \frac{hc}{\lambda} = \frac{(6.626 \times 10^{-34}\text{ J}\cdot\text{s})(3 \times 10^8\text{ m/s})}{632.8 \times 10^{-9}\text{ m}} = 3.141 \times 10^{-19}\text{ J} = 1.96\text{ eV}
\]
The thermal energy at $300\text{ K}$ is $k_B T = (1.38 \times 10^{-23})(300) = 4.14 \times 10^{-21}\text{ J} = 0.02585\text{ eV}$.\\
The ratio is:
\[
\frac{R_{\text{spont}}}{R_{\text{stim}}} = e^{h\nu / k_B T} - 1 = e^{1.96 / 0.02585} - 1 \approx e^{75.8} \approx 8.3 \times 10^{32}
\]
Spontaneous emission is overwhelmingly dominant under thermal equilibrium.
\end{examplebox}

\begin{examplebox}{Worked Example 2.2: Threshold Gain Coefficient Calculation}
\textbf{Problem:} A Ruby laser rod of length $L = 10\text{ cm}$ has mirror reflectivities $R_1 = 1.0$ and $R_2 = 0.90$. The internal scattering loss coefficient is $\alpha_s = 0.02\text{ cm}^{-1}$. Calculate the threshold gain coefficient $\gamma_{\text{th}}$.\\
\textbf{Solution:}\\
Using the threshold gain formula:
\[
\gamma_{\text{th}} = \alpha_s + \frac{1}{2L}\ln\left(\frac{1}{R_1 R_2}\right) = 0.02 + \frac{1}{2(10\text{ cm})}\ln\left(\frac{1}{1.0 \times 0.90}\right)
\]
\[
\gamma_{\text{th}} = 0.02 + \frac{1}{20}\ln(1.1111) = 0.02 + \frac{0.10536}{20} = 0.02 + 0.00527 = 0.02527\text{ cm}^{-1}
\]
The active medium must provide a minimum small-signal optical gain of $0.0253\text{ cm}^{-1}$ to sustain lasing.
\end{examplebox}

\section{Review Questions}
\begin{enumerate}[leftmargin=*,itemsep=3pt]
    \item Derive Einstein's $A$ and $B$ coefficients and explain why lasers cannot operate in thermal equilibrium.
    \item Explain the working mechanism of a 4-level He-Ne laser with an energy level diagram. Why is He gas necessary?
    \item What is population inversion? Derive the expression for the threshold gain coefficient in an optical resonator.
    \item Explain the two-step process of recording and reconstructing a 3D hologram.
\end{enumerate}
